\documentclass[10pt,a4paper]{amsart}
\usepackage{amsmath,amssymb,amsthm,graphicx}
\usepackage{verbatim}
\newenvironment{reference}{\comment}{\endcomment}
%\newenvironment{reference}{}{}
\newenvironment{slogan}{\comment}{\endcomment}
\newenvironment{history}{\comment}{\endcomment}

% For commutative diagrams we use Xy-pic
\usepackage[all]{xy}

% We use 2cell for 2-commutative diagrams.
\xyoption{2cell}
% Theorem environments.
%
\theoremstyle{plain}
\newtheorem{theorem}[subsection]{Theorem}
\newtheorem{proposition}[subsection]{Proposition}
\newtheorem{lemma}[subsection]{Lemma}

\theoremstyle{definition}
\newtheorem{definition}[subsection]{Definition}
\newtheorem{example}[subsection]{Example}
\newtheorem{exercise}[subsection]{Exercise}
\newtheorem{situation}[subsection]{Situation}

\theoremstyle{remark}
\newtheorem{remark}[subsection]{Remark}
\newtheorem{remarks}[subsection]{Remarks}

\numberwithin{equation}{subsection}

% Macros
%
\def\lim{\mathop{\mathrm{lim}}\nolimits}
\def\colim{\mathop{\mathrm{colim}}\nolimits}
\def\Spec{\mathop{\mathrm{Spec}}}
\def\Hom{\mathop{\mathrm{Hom}}\nolimits}
\def\Ext{\mathop{\mathrm{Ext}}\nolimits}
\def\SheafHom{\mathop{\mathcal{H}\!\mathit{om}}\nolimits}
\def\SheafExt{\mathop{\mathcal{E}\!\mathit{xt}}\nolimits}
\def\Sch{\mathit{Sch}}
\def\Mor{\mathop{\mathrm{Mor}}\nolimits}
\def\Ob{\mathop{\mathrm{Ob}}\nolimits}
\def\Sh{\mathop{\mathit{Sh}}\nolimits}
\def\NL{\mathop{N\!L}\nolimits}
\def\CH{\mathop{\mathrm{CH}}\nolimits}
\def\proetale{{pro\text{-}\acute{e}tale}}
\def\etale{{\acute{e}tale}}
\def\QCoh{\mathit{QCoh}}
\def\Ker{\mathop{\mathrm{Ker}}}
\def\Im{\mathop{\mathrm{Im}}}
\def\Coker{\mathop{\mathrm{Coker}}}
\def\Coim{\mathop{\mathrm{Coim}}}

% Boxtimes
%
\DeclareMathSymbol{\boxtimes}{\mathbin}{AMSa}{"02}

%
% Macros for moduli stacks/spaces
%
\def\QCohstack{\mathcal{QC}\!\mathit{oh}}
\def\Cohstack{\mathcal{C}\!\mathit{oh}}
\def\Spacesstack{\mathcal{S}\!\mathit{paces}}
\def\Quotfunctor{\mathrm{Quot}}
\def\Hilbfunctor{\mathrm{Hilb}}
\def\Curvesstack{\mathcal{C}\!\mathit{urves}}
\def\Polarizedstack{\mathcal{P}\!\mathit{olarized}}
\def\Complexesstack{\mathcal{C}\!\mathit{omplexes}}
% \Pic is the operator that assigns to X its picard group, usage \Pic(X)
% \Picardstack_{X/B} denotes the Picard stack of X over B
% \Picardfunctor_{X/B} denotes the Picard functor of X over B
\def\Pic{\mathop{\mathrm{Pic}}\nolimits}
\def\Picardstack{\mathcal{P}\!\mathit{ic}}
\def\Picardfunctor{\mathrm{Pic}}
\def\Deformationcategory{\mathcal{D}\!\mathit{ef}}
\usepackage{cleveref}
\crefname{originallemma}{original lemma}{original lemmas}
\Crefname{originallemma}{Original lemma}{Original lemmas}
\makeatletter
\newlabel{lemma-Noether-normalization}{{00OY}{source}{}{sourcepointer-1}{}}
\newlabel{lemma-Noether-normalization@cref}{{[originallemma][0][]00OY}{[1][0][]source}{}{}{}}
\makeatother
\crefname{originallemma}{original lemma}{original lemmas}
\Crefname{originallemma}{Original lemma}{Original lemmas}
\makeatletter
\newlabel{lemma-quasi-finite-open}{{00QA}{source}{}{sourcepointer-2}{}}
\newlabel{lemma-quasi-finite-open@cref}{{[originallemma][0][]00QA}{[1][0][]source}{}{}{}}
\makeatother
\crefname{originallemma}{original lemma}{original lemmas}
\Crefname{originallemma}{Original lemma}{Original lemmas}
\makeatletter
\newlabel{remark-fundamental-diagram}{{00E6}{source}{}{sourcepointer-3}{}}
\newlabel{remark-fundamental-diagram@cref}{{[originallemma][0][]00E6}{[1][0][]source}{}{}{}}
\makeatother
\title{A finite type algebra of relative dimension n is quasi-finite over a polynomial algebra in n variables after inverting one element}
\author{The Stacks Project Authors}
\date{stacks-project:algebra.tex:local-snapshot-20261012; GNU FDL 1.2 or later}
\begin{document}
\maketitle
\noindent\textit{Editorial note.} All delivered mathematical text is the author's own, copied from the pinned Stacks Project chapter file algebra.tex; no statement or proof step was rewritten, reordered or removed.\par\medskip
\noindent\textit{Editorial note.} Preamble provenance, as required for this source: the frozen snapshot stages the chapter's own preamble file beside it (preamble.tex, 217 lines, 8116 bytes, the same shared Stacks-wide preamble the published chow and etale-cohomology cards use); the first line of the pinned chapter file is \textbackslash{}input\{preamble\}, so no preamble text lives inside the chapter itself. The card takes that file's lines 8-18 (the verbatim package and the three \textbackslash{}comment-based environments reference/slogan/history at lines 9-12, section 13.2, together with the Xy-pic setup) and its lines 154 to the runtime end of the file (the native theorem declarations and mathematical macros). Not one character was changed, no package was added, no macro was defined, and the excerpt serves only as preamble, never as a span coordinate.\par\medskip
\noindent\textit{Editorial note.} The same-chapter pointers are located at their \textbackslash{}label lines inside the pinned chapter file; the cross-chapter pointers are located in the project's own tag list staged in the snapshot (sources/stacks-chow/tags.txt), at the verbatim lines carrying their canonical tags.\par\medskip
\noindent\textit{Editorial note.} The counted proof refers to same-chapter statements by label; they are recorded as navigation stubs pointed at their \textbackslash{}label lines, and the arguments consuming them stay inside the counted span.\par\medskip
\noindent\textit{Editorial note.} The counted spans contain no \textbackslash{}cite, so the card carries no bibliography.\par\medskip
\noindent\textit{Allowed standard prerequisites (editorial).}\begin{itemize}
\item The fibre ring as a finite type algebra over the residue field. Used at the start of the counted proof (lines 31397-31400). Turns the relative dimension hypothesis into a statement about the fibre, where Noether normalisation applies.
\item Quasi-finiteness is an open condition. Invoked at the end of the counted proof (lines 31422-31423). Spreads the quasi-finiteness obtained at one point out to a basic open.
\end{itemize}\medskip
\ifdefined\hypertarget\hypertarget{sourcepointer-1}{}\fi
\noindent\textit{Source reference.} Original lemma 00OY: sources/stacks-algebra/pinned-original-source.tex, source lines 28242--28242. Navigation pointer for \textbackslash{}ref\{lemma-Noether-normalization\} used inside the counted span (line 31416); the argument that consumes it is retained in the span.\par\medskip
\ifdefined\hypertarget\hypertarget{sourcepointer-2}{}\fi
\noindent\textit{Source reference.} Original lemma 00QA: sources/stacks-algebra/pinned-original-source.tex, source lines 31097--31097. Navigation pointer for \textbackslash{}ref\{lemma-quasi-finite-open\} used inside the counted span (line 31422); the argument that consumes it is retained in the span.\par\medskip
\ifdefined\hypertarget\hypertarget{sourcepointer-3}{}\fi
\noindent\textit{Source reference.} Original lemma 00E6: sources/stacks-algebra/pinned-original-source.tex, source lines 3639--3639. Navigation pointer for \textbackslash{}ref\{remark-fundamental-diagram\} used inside the counted span (line 31407); the argument that consumes it is retained in the span.\par\medskip
\begin{lemma}
\label{lemma-quasi-finite-over-polynomial-algebra}
Let $R \to S$ be a finite type ring map.
Let $\mathfrak q \subset S$ be a prime.
Suppose that $\dim_{\mathfrak q}(S/R) = n$.
There exists a $g \in S$, $g \not\in \mathfrak q$
such that $S_g$ is quasi-finite over a
polynomial algebra $R[t_1, \ldots, t_n]$.
\end{lemma}

\begin{proof}
The ring $\overline{S} = S \otimes_R \kappa(\mathfrak p)$ is
of finite type over $\kappa(\mathfrak p)$.
Let $\overline{\mathfrak q}$ be the prime of $\overline{S}$
corresponding to $\mathfrak q$.
By definition of
the dimension of a topological space at a point there exists
an open $U \subset \Spec(\overline{S})$ with
$\overline{q} \in U$ and $\dim(U) = n$.
Since the topology on $\Spec(\overline{S})$ is
induced from the topology on $\Spec(S)$ (see
Remark \ref{remark-fundamental-diagram}), we can find
a $g \in S$, $g \not \in \mathfrak q$ with image
$\overline{g} \in \overline{S}$ such that
$D(\overline{g}) \subset U$.
Thus after replacing $S$ by $S_g$ we see that
$\dim(\overline{S}) = n$.

\medskip\noindent
Next, choose generators $x_1, \ldots, x_N$ for $S$ as an $R$-algebra. By
Lemma \ref{lemma-Noether-normalization}
there exist elements $y_1, \ldots, y_n$ in the $\mathbf{Z}$-subalgebra of $S$
generated by $x_1, \ldots, x_N$ such that the map
$R[t_1, \ldots, t_n] \to S$, $t_i \mapsto y_i$ has the property
that $\kappa(\mathfrak p)[t_1\ldots, t_n] \to \overline{S}$
is finite. In particular, $S$ is quasi-finite over $R[t_1, \ldots, t_n]$
at $\mathfrak q$. Hence, by Lemma \ref{lemma-quasi-finite-open}
we may replace $S$ by $S_g$ for some $g\in S$, $g \not \in \mathfrak q$
such that $R[t_1, \ldots, t_n] \to S$ is quasi-finite.
\end{proof}

\end{document}
